
\section{Querverweise}

\label{sec:Gliederung}

Hier füge ich Querverweise ein. Dafür benutze ich ein Label und eine Referenz. 
Das Label ist die Stelle, wo ich hinverweisen will. 
\textbackslash label bezieht sich immer auf den zuletzt nummerierten Gliederungs-/Zählerstand.
Die Referenz ist der Ausgangspunkt.
Dem Label kann ich ein Parameter mitgeben für diese Stelle im Text. \\
Es ist Konvention ein Präfix davor zu setzen. Das könnte man auch weglassen, ist aber
praktisch, um zwischen Bildern, Tabellen und nummerierten Elementen unterscheiden zu 
können, wenn man später darauf referenzieren will. \\
\noindent Mögliche Präfixe: 
\begin{description}
    \item [sec] section
    \item [cha] chapter 
    \item [fig] Abbildungen
    \item [tab] Tabellen
\end{description}
Mit \textbackslash ref verweise ich auf die Stelle, wobei die Nummerierung im pdf automatisch erfolgt. 
Mit \textbackslash pageref verweise ich auf die Seitenzahl der angegebenen Label-Markierung. \\
Mit \textbackslash nameref referenziere ich auf den Namen der Abbildung (o.ä.).
Hier kommen jetzt die Verweise erstens auf dieses Kapitel \ref{sec:Gliederung} und zweitens auf 
die Abbildung \nameref{fig:Raspie} vom Rasperry Pie \ref{fig:Raspie} auf Seite \pageref{sec:Gliederung}. \\
Einfacher kann ich mir das machen, indem ich einen neuen Befehl definiere. So verlinke ich jetzt
mal auf die \Tabelle{Super-Tabelle}. Wenn ich das ganze auf einen Befehl auslagere, bleibt der Text 
hier etwas übersichtlicher. \\
Ebenfalls ist es möglich das ganze in eine Fußnote auszulagern. Das versuchen wir jetzt noch
mal mit einem Verweis auf die verschiedenen Schriftauszeichnungen\footnote{Siehe \Kapitel{Schriftauszeichnung}.}. Hierfür benutzen wir \textbackslash footnote.

